\chapter*{Preface}\label{sec:Preface}
% (1 pagina) \\
% Cuales son los antecedentes que motivan este proyecto: por que este proyecto es relevante en este momento?
% la motivación por que es necesario checar el soc de la moto? \\
% Pues que se estan utilizando vehículos eléctricos ahora. Y por seguridad se tiene que vigilar las baterias

% (maybe i can give more fun facts here (?))\\
The transportation sector is the primary source of direct greenhouse gas emissions in the United States, accounting 
for 28\% of the national total~\cite{epa_transportation_2024}. In the European Union, this sector is responsible for
approximately one-fifth of total \ac{CO2} emissions. The shift towards electric mobility is being accelerated by regulations, 
such as the EU's plan for all new cars and vans to be zero-emission by 2035. Consequently, the EU estimates that battery-electric vehicles
will dominate the market~\cite{europarl_2022_cars}.
\\
\\
\acp{EV} are constrained by their battery energy capacity, as the size and energy density of the battery pack 
determine the vehicle's driving range~\cite{plett2015battery}.~\ac{Li-ion} batteries have played a 
major role in the decarbonisation of mobility due to their lightweight characteristics and high efficiency.
\ac{Li-ion} batteries have addressed these limitations and have emerged as a widespread technology for 
electric mobility~\cite{IPCC2022_Mitigation}, replacing and outperforming other battery technologies such as lead-acid 
or nickel-metal hydride~\cite{ma16176063}. However, their performance
depends on factors such as the temperature, current loads, and charge-discharge cycles~\cite{Bosello2025}. 
Another critical aspect concerns the safety risks associated 
with overcharging or overdischarging, as these can trigger 
internal chemical reactions 
or short-circuiting~\cite{plett2015battery}. Due to these challenges and to guarantee optimal performance, longevity, and safety, 
accurate monitoring of the battery state, as well as individual cell voltages and temperatures, is essential~\cite{Bosello2025}.
\\
\\
% \textit{(here saying that the SOC and SOH is super important, why do we measure them? how do we measure them?)} \\
Among critical internal states required for effective monitoring are the \ac{SOC} and \ac{SOH}. 
As these variables cannot be directly measured by any sensor, they must be estimated from measurable
signals, such as voltage, current, and temperature. Estimation algorithms vary in complexity and accuracy. 
In high-performance and 
safety-critical systems, the consequences of estimation errors extend beyond suboptimal performance and 
may constitute a safety risk. Inaccurate \ac{SOC} estimation can 
lead to an inefficient use of battery capacity and cell damage. Similarly, incorrect \ac{SOH} can lead to unpredictable system 
behaviour and unreliability~\cite{plett2015battery}.
\\
% \\
% \textit{(and then taking all that i have said into the racing motorcycle, like maybe what they are using, or maybe like yeah they are using 
% this type of battery and reach this voltage, and say wow that's a really high voltage, jesus, that would be dangerous if someone is 
% not checking on that. how do WE the race team read the voltage and temp from the battery --> BMS!)}
\\
The eRacing team at ETSEIB designs and develops an electric motorcycle to compete at the international MotoStudent competition.
To meet the performance requirements at different track conditions~\cite{MotoStudent2026Regulations}, the vehicle is powered by \ac{NMC} 
batteries, a type of \ac{Li-ion} battery chemistry. The system operates at high voltages, which requires safety protocols to mitigate operational 
hazards. 
Consequently, the \ac{BMS} is responsible for monitoring individual cell voltages and temperatures in real-time, as a safety requirement and
as a necessary input for the State of Charge and State of Health estimation. However, achieving an accurate real-time state estimation 
remains a challenge, 
particularly under the dynamic current/voltage profiles of a race.
\\
\\
This work addresses this challenge by developing and validating a model-based dual estimation framework for simultaneous 
\ac{SOC} and \ac{SOH} estimation of a $30$ Ah / $111$ Wh \ac{NMC} battery~\cite{welion2026shp270} used in the eRacing electric motorcycle. 
A second-order RC equivalent circuit model is parameterised through experimental characterisation, and dual 
Extended Kalman Filters are implemented to estimate battery states under dynamic operating conditions.
